\begin{table}[t]
\centering
\caption{P3c M2: area\_total 분해 — 비조건부 (Ch4+Ch5) vs 조건부 (Ch4 단독)}
\label{tab:ch5_area_decomposition_ko}
\begin{tabular}{llrrr}
\toprule
Bandwidth & Year & Unconditional & Conditional (n=5) & Selection share \\
\midrule
T1 & 2018 & -190 & -57 & 70.2% \\
T1 & 2020 & -10 & 17 & 275.0% \\
T1 & 2021 & 237 & 311 & -31.1% \\
T1 & 2022 & 23 & 187 & -725.6% \\
T2 & 2018 & -61 & -29 & 53.0% \\
T2 & 2020 & 218 & 217 & 0.1% \\
T2 & 2021 & 317 & 320 & -1.0% \\
T2 & 2022 & 358 & 365 & -2.0% \\
T3 & 2018 & 33 & 90 & -177.6% \\
T3 & 2020 & 176* & 174 & 1.1% \\
T3 & 2021 & 392*** & 377** & 3.8% \\
T3 & 2022 & 493*** & 497*** & -0.9% \\
\bottomrule
\end{tabular}\\
\footnotesize\textit{비조건부 컬럼 = ch3\_area\_events (06\_channels.R Phase 2, 전체 패널). 조건부 컬럼 = n\_years==5 완전 농가만 (N=2,217). 선택 비중 = 1 - (조건부 / 비조건부). Channel 4(iii) 잔존 농가 내 확장 마진 = 조건부; Channel 5 선택 편의 = 비조건부 - 조건부 비율. * p$<$0.10, ** p$<$0.05, *** p$<$0.01.}
\end{table}

